\IfFileExists{stacks-project.cls}{%
\documentclass{stacks-project}
}{%
\documentclass{amsart}
}

% For dealing with references we use the comment environment
\usepackage{verbatim}
\newenvironment{reference}{\comment}{\endcomment}
%\newenvironment{reference}{}{}
\newenvironment{slogan}{\comment}{\endcomment}
\newenvironment{history}{\comment}{\endcomment}

% For commutative diagrams we use Xy-pic
\usepackage[all]{xy}

% We use 2cell for 2-commutative diagrams.
\xyoption{2cell}
\UseAllTwocells

% We use multicol for the list of chapters between chapters
\usepackage{multicol}

% This is generally recommended for better output
\usepackage{lmodern}
\usepackage[T1]{fontenc}

% For cross-file-references
\usepackage{xr-hyper}

% Package for hypertext links:
\usepackage{hyperref}

% For any local file, say "hello.tex" you want to link to please
% use \externaldocument[hello-]{hello}
\externaldocument[introduction-]{introduction}
\externaldocument[conventions-]{conventions}
\externaldocument[sets-]{sets}
\externaldocument[categories-]{categories}
\externaldocument[topology-]{topology}
\externaldocument[sheaves-]{sheaves}
\externaldocument[sites-]{sites}
\externaldocument[stacks-]{stacks}
\externaldocument[fields-]{fields}
\externaldocument[algebra-]{algebra}
\externaldocument[brauer-]{brauer}
\externaldocument[homology-]{homology}
\externaldocument[derived-]{derived}
\externaldocument[simplicial-]{simplicial}
\externaldocument[more-algebra-]{more-algebra}
\externaldocument[smoothing-]{smoothing}
\externaldocument[modules-]{modules}
\externaldocument[sites-modules-]{sites-modules}
\externaldocument[injectives-]{injectives}
\externaldocument[cohomology-]{cohomology}
\externaldocument[sites-cohomology-]{sites-cohomology}
\externaldocument[dga-]{dga}
\externaldocument[dpa-]{dpa}
\externaldocument[sdga-]{sdga}
\externaldocument[hypercovering-]{hypercovering}
\externaldocument[schemes-]{schemes}
\externaldocument[constructions-]{constructions}
\externaldocument[properties-]{properties}
\externaldocument[morphisms-]{morphisms}
\externaldocument[coherent-]{coherent}
\externaldocument[divisors-]{divisors}
\externaldocument[limits-]{limits}
\externaldocument[varieties-]{varieties}
\externaldocument[topologies-]{topologies}
\externaldocument[descent-]{descent}
\externaldocument[perfect-]{perfect}
\externaldocument[more-morphisms-]{more-morphisms}
\externaldocument[flat-]{flat}
\externaldocument[groupoids-]{groupoids}
\externaldocument[more-groupoids-]{more-groupoids}
\externaldocument[etale-]{etale}
\externaldocument[chow-]{chow}
\externaldocument[intersection-]{intersection}
\externaldocument[pic-]{pic}
\externaldocument[weil-]{weil}
\externaldocument[adequate-]{adequate}
\externaldocument[dualizing-]{dualizing}
\externaldocument[duality-]{duality}
\externaldocument[discriminant-]{discriminant}
\externaldocument[derham-]{derham}
\externaldocument[local-cohomology-]{local-cohomology}
\externaldocument[algebraization-]{algebraization}
\externaldocument[curves-]{curves}
\externaldocument[resolve-]{resolve}
\externaldocument[models-]{models}
\externaldocument[functors-]{functors}
\externaldocument[equiv-]{equiv}
\externaldocument[pione-]{pione}
\externaldocument[etale-cohomology-]{etale-cohomology}
\externaldocument[proetale-]{proetale}
\externaldocument[relative-cycles-]{relative-cycles}
\externaldocument[more-etale-]{more-etale}
\externaldocument[trace-]{trace}
\externaldocument[crystalline-]{crystalline}
\externaldocument[spaces-]{spaces}
\externaldocument[spaces-properties-]{spaces-properties}
\externaldocument[spaces-morphisms-]{spaces-morphisms}
\externaldocument[decent-spaces-]{decent-spaces}
\externaldocument[spaces-cohomology-]{spaces-cohomology}
\externaldocument[spaces-limits-]{spaces-limits}
\externaldocument[spaces-divisors-]{spaces-divisors}
\externaldocument[spaces-over-fields-]{spaces-over-fields}
\externaldocument[spaces-topologies-]{spaces-topologies}
\externaldocument[spaces-descent-]{spaces-descent}
\externaldocument[spaces-perfect-]{spaces-perfect}
\externaldocument[spaces-more-morphisms-]{spaces-more-morphisms}
\externaldocument[spaces-flat-]{spaces-flat}
\externaldocument[spaces-groupoids-]{spaces-groupoids}
\externaldocument[spaces-more-groupoids-]{spaces-more-groupoids}
\externaldocument[bootstrap-]{bootstrap}
\externaldocument[spaces-pushouts-]{spaces-pushouts}
\externaldocument[spaces-chow-]{spaces-chow}
\externaldocument[groupoids-quotients-]{groupoids-quotients}
\externaldocument[spaces-more-cohomology-]{spaces-more-cohomology}
\externaldocument[spaces-simplicial-]{spaces-simplicial}
\externaldocument[spaces-duality-]{spaces-duality}
\externaldocument[formal-spaces-]{formal-spaces}
\externaldocument[restricted-]{restricted}
\externaldocument[spaces-resolve-]{spaces-resolve}
\externaldocument[formal-defos-]{formal-defos}
\externaldocument[defos-]{defos}
\externaldocument[cotangent-]{cotangent}
\externaldocument[examples-defos-]{examples-defos}
\externaldocument[algebraic-]{algebraic}
\externaldocument[examples-stacks-]{examples-stacks}
\externaldocument[stacks-sheaves-]{stacks-sheaves}
\externaldocument[criteria-]{criteria}
\externaldocument[artin-]{artin}
\externaldocument[quot-]{quot}
\externaldocument[stacks-properties-]{stacks-properties}
\externaldocument[stacks-morphisms-]{stacks-morphisms}
\externaldocument[stacks-limits-]{stacks-limits}
\externaldocument[stacks-cohomology-]{stacks-cohomology}
\externaldocument[stacks-perfect-]{stacks-perfect}
\externaldocument[stacks-introduction-]{stacks-introduction}
\externaldocument[stacks-more-morphisms-]{stacks-more-morphisms}
\externaldocument[stacks-geometry-]{stacks-geometry}
\externaldocument[moduli-]{moduli}
\externaldocument[moduli-curves-]{moduli-curves}
\externaldocument[examples-]{examples}
\externaldocument[exercises-]{exercises}
\externaldocument[guide-]{guide}
\externaldocument[desirables-]{desirables}
\externaldocument[coding-]{coding}
\externaldocument[obsolete-]{obsolete}
\externaldocument[fdl-]{fdl}
\externaldocument[index-]{index}

% Theorem environments.
%
\theoremstyle{plain}
\newtheorem{theorem}[subsection]{Theorem}
\newtheorem{proposition}[subsection]{Proposition}
\newtheorem{lemma}[subsection]{Lemma}

\theoremstyle{definition}
\newtheorem{definition}[subsection]{Definition}
\newtheorem{example}[subsection]{Example}
\newtheorem{exercise}[subsection]{Exercise}
\newtheorem{situation}[subsection]{Situation}

\theoremstyle{remark}
\newtheorem{remark}[subsection]{Remark}
\newtheorem{remarks}[subsection]{Remarks}

\numberwithin{equation}{subsection}

% Macros
%
\def\lim{\mathop{\mathrm{lim}}\nolimits}
\def\colim{\mathop{\mathrm{colim}}\nolimits}
\def\Spec{\mathop{\mathrm{Spec}}}
\def\Hom{\mathop{\mathrm{Hom}}\nolimits}
\def\Ext{\mathop{\mathrm{Ext}}\nolimits}
\def\SheafHom{\mathop{\mathcal{H}\!\mathit{om}}\nolimits}
\def\SheafExt{\mathop{\mathcal{E}\!\mathit{xt}}\nolimits}
\def\Sch{\mathit{Sch}}
\def\Mor{\mathop{\mathrm{Mor}}\nolimits}
\def\Ob{\mathop{\mathrm{Ob}}\nolimits}
\def\Sh{\mathop{\mathit{Sh}}\nolimits}
\def\NL{\mathop{N\!L}\nolimits}
\def\CH{\mathop{\mathrm{CH}}\nolimits}
\def\proetale{{pro\text{-}\acute{e}tale}}
\def\etale{{\acute{e}tale}}
\def\QCoh{\mathit{QCoh}}
\def\Ker{\mathop{\mathrm{Ker}}}
\def\Im{\mathop{\mathrm{Im}}}
\def\Coker{\mathop{\mathrm{Coker}}}
\def\Coim{\mathop{\mathrm{Coim}}}

% Boxtimes
%
\DeclareMathSymbol{\boxtimes}{\mathbin}{AMSa}{"02}

%
% Macros for moduli stacks/spaces
%
\def\QCohstack{\mathcal{QC}\!\mathit{oh}}
\def\Cohstack{\mathcal{C}\!\mathit{oh}}
\def\Spacesstack{\mathcal{S}\!\mathit{paces}}
\def\Quotfunctor{\mathrm{Quot}}
\def\Hilbfunctor{\mathrm{Hilb}}
\def\Curvesstack{\mathcal{C}\!\mathit{urves}}
\def\Polarizedstack{\mathcal{P}\!\mathit{olarized}}
\def\Complexesstack{\mathcal{C}\!\mathit{omplexes}}
% \Pic is the operator that assigns to X its picard group, usage \Pic(X)
% \Picardstack_{X/B} denotes the Picard stack of X over B
% \Picardfunctor_{X/B} denotes the Picard functor of X over B
\def\Pic{\mathop{\mathrm{Pic}}\nolimits}
\def\Picardstack{\mathcal{P}\!\mathit{ic}}
\def\Picardfunctor{\mathrm{Pic}}
\def\Deformationcategory{\mathcal{D}\!\mathit{ef}}
\usepackage{cleveref}
\crefname{originallemma}{original lemma}{original lemmas}
\Crefname{originallemma}{Original lemma}{Original lemmas}
\makeatletter
\newlabel{lemma-quot-sheaf}{{09TT}{source}{}{sourcepointer-1}{}}
\newlabel{lemma-quot-sheaf@cref}{{[originallemma][0][]09TT}{[1][0][]source}{}{}{}}
\makeatother
\crefname{originalsection}{original section}{original sections}
\Crefname{originalsection}{Original section}{Original sections}
\makeatletter
\newlabel{algebraic-section-split}{{04SU}{source}{}{sourcepointer-2}{}}
\newlabel{algebraic-section-split@cref}{{[originalsection][0][]04SU}{[1][0][]source}{}{}{}}
\makeatother
\crefname{originalproposition}{original proposition}{original propositions}
\Crefname{originalproposition}{Original proposition}{Original propositions}
\makeatletter
\newlabel{algebraic-proposition-algebraic-stack-no-automorphisms}{{04SZ}{source}{}{sourcepointer-3}{}}
\newlabel{algebraic-proposition-algebraic-stack-no-automorphisms@cref}{{[originalproposition][0][]04SZ}{[1][0][]source}{}{}{}}
\makeatother
\crefname{originaltheorem}{original theorem}{original theorems}
\Crefname{originaltheorem}{Original theorem}{Original theorems}
\makeatletter
\newlabel{theorem-coherent-algebraic-general}{{09DS}{source}{}{sourcepointer-4}{}}
\newlabel{theorem-coherent-algebraic-general@cref}{{[originaltheorem][0][]09DS}{[1][0][]source}{}{}{}}
\makeatother
\crefname{originallemma}{original lemma}{original lemmas}
\Crefname{originallemma}{Original lemma}{Original lemmas}
\makeatletter
\newlabel{algebraic-lemma-representable-morphism-to-algebraic}{{05UM}{source}{}{sourcepointer-5}{}}
\newlabel{algebraic-lemma-representable-morphism-to-algebraic@cref}{{[originallemma][0][]05UM}{[1][0][]source}{}{}{}}
\makeatother
\crefname{originallemma}{original lemma}{original lemmas}
\Crefname{originallemma}{Original lemma}{Original lemmas}
\makeatletter
\newlabel{stacks-lemma-2-fibre-product-gives-stack-in-setoids}{{05UI}{source}{}{sourcepointer-6}{}}
\newlabel{stacks-lemma-2-fibre-product-gives-stack-in-setoids@cref}{{[originallemma][0][]05UI}{[1][0][]source}{}{}{}}
\makeatother
\crefname{originallemma}{original lemma}{original lemmas}
\Crefname{originallemma}{Original lemma}{Original lemmas}
\makeatletter
\newlabel{stacks-lemma-when-stack-in-sets}{{0430}{source}{}{sourcepointer-7}{}}
\newlabel{stacks-lemma-when-stack-in-sets@cref}{{[originallemma][0][]0430}{[1][0][]source}{}{}{}}
\makeatother
\crefname{originallemma}{original lemma}{original lemmas}
\Crefname{originallemma}{Original lemma}{Original lemmas}
\makeatletter
\newlabel{spaces-flat-lemma-F-surj-open}{{09TP}{source}{}{sourcepointer-8}{}}
\newlabel{spaces-flat-lemma-F-surj-open@cref}{{[originallemma][0][]09TP}{[1][0][]source}{}{}{}}
\makeatother
\crefname{originalproposition}{original proposition}{original propositions}
\Crefname{originalproposition}{Original proposition}{Original propositions}
\makeatletter
\newlabel{proposition-hom}{{08K6}{source}{}{sourcepointer-9}{}}
\newlabel{proposition-hom@cref}{{[originalproposition][0][]08K6}{[1][0][]source}{}{}{}}
\makeatother
\title{The Quot functor is an algebraic space, locally of finite presentation}
\author{The Stacks Project Authors}
\date{stacks-project:quot.tex:local-snapshot-20261012; GNU FDL 1.2 or later}
\begin{document}
\maketitle
\setcounter{subsection}{1}
\noindent\textit{Editorial note.} All delivered mathematical text is the author's own, copied from the pinned Stacks Project chapter file quot.tex; no statement or proof step was rewritten, reordered or removed. The counted span is the author environment and its own proof environment. References inside the span are handed to the assembler as navigation stubs that carry the canonical tag of the cited tag only; whether each of them is load-bearing is recorded on the stub itself.\par\medskip
\noindent\textit{Allowed standard prerequisites (editorial).}\begin{itemize}
\item Lemma lemma-quot-sheaf (canonical tag 09TT). cited at proof line 2540. The Quot functor is an fppf sheaf.
\item Section algebraic-section-split (canonical tag 04SU). cited at proof line 2544. Passing from the sheaf to the stack in groupoids.
\item Proposition algebraic-proposition-algebraic-stack-no-automorphisms (canonical tag 04SZ). cited at proof lines 2545-2546. Reduction to algebraicity of the Quot stack.
\item Theorem theorem-coherent-algebraic-general (canonical tag 09DS). cited at proof line 2557. Cohstack\_\{X/B\} is an algebraic stack.
\item Lemma algebraic-lemma-representable-morphism-to-algebraic (canonical tag 05UM). cited at proof lines 2559-2560. Suffices to show the 1-morphism is representable by algebraic spaces.
\item Lemma stacks-lemma-2-fibre-product-gives-stack-in-setoids (canonical tag 05UI). cited at proof lines 2576 and 2579. Z is a stack in setoids; its sheaf of sets.
\item Lemma stacks-lemma-when-stack-in-sets (canonical tag 0430). cited at proof line 2580. Identifying the sheaf of sets attached to Z.
\item Lemma spaces-flat-lemma-F-surj-open (canonical tag 09TP). cited at proof line 2584. Surjectivity locus is open: Z is an open subspace.
\item Proposition proposition-hom (canonical tag 08K6). cited at proof line 2587. The ambient algebraic space Hom(F\_T, G).
\end{itemize}\medskip
\ifdefined\hypertarget\hypertarget{sourcepointer-1}{}\fi
\noindent\textit{Source reference.} Original lemma 09TT: sources/stacks-chow/tags.txt, source lines 12401--12401. Cited at proof line 2540: sheaf property of the Quot functor; load-bearing.\par\medskip
\ifdefined\hypertarget\hypertarget{sourcepointer-2}{}\fi
\noindent\textit{Source reference.} Original section 04SU: sources/stacks-chow/tags.txt, source lines 6232--6232. Cited at proof line 2544: passing from the sheaf to the associated stack in groupoids; load-bearing.\par\medskip
\ifdefined\hypertarget\hypertarget{sourcepointer-3}{}\fi
\noindent\textit{Source reference.} Original proposition 04SZ: sources/stacks-chow/tags.txt, source lines 6237--6237. Cited at proof lines 2545-2546: reduces to showing the Quot stack is algebraic; load-bearing.\par\medskip
\ifdefined\hypertarget\hypertarget{sourcepointer-4}{}\fi
\noindent\textit{Source reference.} Original theorem 09DS: sources/stacks-chow/tags.txt, source lines 11875--11875. Cited at proof line 2557: Cohstack\_\{X/B\} is an algebraic stack; load-bearing.\par\medskip
\ifdefined\hypertarget\hypertarget{sourcepointer-5}{}\fi
\noindent\textit{Source reference.} Original lemma 05UM: sources/stacks-chow/tags.txt, source lines 7521--7521. Cited at proof lines 2559-2560: suffices to show the 1-morphism is representable by algebraic spaces; load-bearing.\par\medskip
\ifdefined\hypertarget\hypertarget{sourcepointer-6}{}\fi
\noindent\textit{Source reference.} Original lemma 05UI: sources/stacks-chow/tags.txt, source lines 7517--7517. Cited at proof lines 2576 and 2579: the fibre product Z is a stack in setoids and corresponds to a sheaf; load-bearing twice.\par\medskip
\ifdefined\hypertarget\hypertarget{sourcepointer-7}{}\fi
\noindent\textit{Source reference.} Original lemma 0430: sources/stacks-chow/tags.txt, source lines 5347--5347. Cited at proof line 2580: identifying the sheaf of sets attached to Z; load-bearing.\par\medskip
\ifdefined\hypertarget\hypertarget{sourcepointer-8}{}\fi
\noindent\textit{Source reference.} Original lemma 09TP: sources/stacks-chow/tags.txt, source lines 12397--12397. Cited at proof line 2584: Z is representable by an open subspace of the Hom space; load-bearing.\par\medskip
\ifdefined\hypertarget\hypertarget{sourcepointer-9}{}\fi
\noindent\textit{Source reference.} Original proposition 08K6: sources/stacks-chow/tags.txt, source lines 10874--10874. Cited at proof line 2587: the ambient algebraic space Hom(F\_T, G); load-bearing.\par\medskip
\begin{proposition}
\label{proposition-quot}
Let $S$ be a scheme. Let $f : X \to B$ be a morphism of algebraic
spaces over $S$. Let $\mathcal{F}$ be a quasi-coherent sheaf
on $X$. If $f$ is of finite presentation and separated, then
$\Quotfunctor_{\mathcal{F}/X/B}$
is an algebraic space. If $\mathcal{F}$ is of finite presentation,
then $\Quotfunctor_{\mathcal{F}/X/B} \to B$ is locally of finite presentation.
\end{proposition}

\begin{proof}
By Lemma \ref{lemma-quot-sheaf}
we have that $\Quotfunctor_{\mathcal{F}/X/B}$ is a sheaf in the
fppf topology. Let $\textit{Quot}_{\mathcal{F}/X/B}$ be the stack in
groupoids corresponding to $\Quotfunctor_{\mathcal{F}/X/S}$, see
Algebraic Stacks, Section \ref{algebraic-section-split}.
By Algebraic Stacks, Proposition
\ref{algebraic-proposition-algebraic-stack-no-automorphisms}
it suffices to show that $\textit{Quot}_{\mathcal{F}/X/B}$
is an algebraic stack.
Consider the $1$-morphism of stacks in groupoids
$$
\textit{Quot}_{\mathcal{F}/X/S}
\longrightarrow
\Cohstack_{X/B}
$$
on $(\Sch/S)_{fppf}$ which associates to the quotient
$\mathcal{F}_T \to \mathcal{Q}$ the module $\mathcal{Q}$.
By Theorem \ref{theorem-coherent-algebraic-general} we know that
$\Cohstack_{X/B}$ is an algebraic stack.
By Algebraic Stacks, Lemma
\ref{algebraic-lemma-representable-morphism-to-algebraic}
it suffices to show that this $1$-morphism is representable
by algebraic spaces.

\medskip\noindent
Let $T$ be a scheme over $S$ and let the object $(h, \mathcal{G})$ of
$\Cohstack_{X/B}$ over $T$ correspond
to a $1$-morphism $\xi : (\Sch/T)_{fppf} \to \Cohstack_{X/B}$.
The $2$-fibre product
$$
\mathcal{Z} =
(\Sch/T)_{fppf}
\times_{\xi, \Cohstack_{X/B}}
\textit{Quot}_{\mathcal{F}/X/S}
$$
is a stack in setoids, see
Stacks, Lemma \ref{stacks-lemma-2-fibre-product-gives-stack-in-setoids}.
The corresponding sheaf of sets (i.e., functor, see
Stacks, Lemmas
\ref{stacks-lemma-2-fibre-product-gives-stack-in-setoids} and
\ref{stacks-lemma-when-stack-in-sets}) assigns to a scheme
$T'/T$ the set of surjections $u : \mathcal{F}_{T'} \to \mathcal{G}_{T'}$
of quasi-coherent modules on $X_{T'}$. Thus we see that
$\mathcal{Z}$ is representable by an open subspace
(by Flatness on Spaces, Lemma \ref{spaces-flat-lemma-F-surj-open})
of the algebraic space
$\mathit{Hom}(\mathcal{F}_T, \mathcal{G})$ from
Proposition \ref{proposition-hom}.
\end{proof}

\end{document}
